\documentclass[11pt]{amsart}
\usepackage{amsmath,amssymb,amsthm,mathtools}
\usepackage[margin=1in]{geometry}
\usepackage{booktabs}
\usepackage{enumitem}
\usepackage[colorlinks=true,linkcolor=blue,citecolor=blue]{hyperref}

\theoremstyle{plain}
\newtheorem{theorem}{Theorem}[section]
\newtheorem{lemma}[theorem]{Lemma}
\newtheorem{corollary}[theorem]{Corollary}
\newtheorem{proposition}[theorem]{Proposition}
\newtheorem{conjecture}[theorem]{Conjecture}

\theoremstyle{definition}
\newtheorem{definition}[theorem]{Definition}

\theoremstyle{remark}
\newtheorem*{remark}{Remark}
\newtheorem*{observation}{Observation}

\newcommand{\Gs}{G^{*}}
\newcommand{\Gss}{{G^{*}}^{2}}
\newcommand{\Gsss}{{G^{*}}^{3}}
\newcommand{\GG}{G_{\!\scriptscriptstyle\mathrm{G}}}
\newcommand{\Z}{\mathbb{Z}}
\newcommand{\Q}{\mathbb{Q}}
\newcommand{\R}{\mathbb{R}}
\newcommand{\C}{\mathbb{C}}
\newcommand{\Aut}{\mathrm{Aut}}

\title{The master quadratic $P_{\Gs}$ as a math--physics bridge:
$\chi_{-4}$, the ternary substrate, and the FTD identification}
\author{}
\date{2026-05-19}

\begin{document}
\maketitle

\begin{abstract}
The companion paper \cite{paper-gstar-A} establishes that the Kronecker
character $\chi_{-4}$ of $K = \Q(i)$ generates the entire algebraic--analytic
identity ring of the lemniscatic constant $\Gs = \Gamma(1/4)/\Gamma(3/4)$
through four functorial projections (lattice, Chowla--Selberg, Hecke,
Dirichlet), and that $\chi_{-4}(n) = \mathrm{Im}(i^{n})$ with value set
$\{-1, 0, +1\}$. This companion note draws the consequences of those
results for the math--physics bridge of the Foundational Ternary Dynamics
(FTD) framework \cite{ftd-spec}, where the same set $\{-1, 0, +1\}$
appears as the ternary voxel alphabet of a discrete physical substrate.
Three structural observations are recorded, all explicitly tagged as
synthesis rather than theorem: (i) the FTD voxel alphabet coincides as
a set with the value set of $\chi_{-4}$, sourcing a candidate
\emph{Gaussian-integer-module ontology}; (ii) the Born rule
$\psi \mapsto |\psi|^{2}$ and the character $\chi_{-4}: i^{n} \mapsto
\mathrm{Im}(i^{n})$ are the same arrow --- a real-coordinate projection
of a complex object --- applied at different scopes; (iii) the FTD
identification $(x_{+}, x_{-}) \approx (\alpha^{-1}, N_{c})$ of the
roots of the master quadratic $P_{\Gs}(x) = x^{2} - 16\Gss x + 16\Gsss$
with the fine-structure constant and the number of QCD colours requires
exactly one bit of physical input --- the exponent-pair selector $(2, 3)$
--- which is over-determined by joint root-matching. All mathematical
content is taken from \cite{paper-gstar-A}; this companion contributes
only the bridge interpretation, presented under the project's
\texttt{[SYNTHESIS]} tag.
\end{abstract}

\section{Setting and scope}\label{sec:setting}

We adopt the entire mathematical content of \cite{paper-gstar-A}
without re-derivation. The constants in use are:
\begin{itemize}[leftmargin=*, itemsep=2pt]
\item $\Gs = \Gamma(1/4)/\Gamma(3/4) \approx 2.9587$;
\item $\GG = 1/\mathrm{AGM}(1, \sqrt{2}) \approx 0.8346$, with the bridge identity $\Gs = 2\sqrt{\pi}\,\GG$;
\item $\chi_{-4}: (\Z/4\Z)^\times \to \{\pm 1\}$, the unique non-trivial Dirichlet character mod $4$;
\item the master quadratic $P_{\Gs}(x) = x^{2} - 16\,\Gss\,x + 16\,\Gsss$, with integer coefficient $16 = |\Aut_\C(E_{\mathrm{lemn}})|^{2}$ and roots $x_{\pm}$ satisfying Vieta's identities $x_+ + x_- = 16\Gss$, $x_+ x_- = 16\Gsss$, $1/x_+ + 1/x_- = 1/\Gs$.
\end{itemize}
The mathematical paper \cite{paper-gstar-A} proves that $\chi_{-4}$
admits the explicit form $\chi_{-4}(n) = \mathrm{Im}(i^{n}) = \sin(\pi n/2)$
and that its value set is
\[
\{\chi_{-4}(n) : n \in \Z\} \;=\; \{i^{2},\, 0,\, |i^{2}|\} \;=\; \{-1, 0, +1\}.
\]
Numerical observation: the roots of $P_{\Gs}$ are
$x_{+} \approx 137.036\ldots$ (matching the inverse fine-structure constant
$\alpha^{-1}$ to 1.26~ppm) and $x_{-} \approx 3.024\ldots$ (matching
the number of QCD colours $N_{c} = 3$ to 0.8\%). The Foundational
Ternary Dynamics framework \cite{ftd-spec} posits the identification
$(x_{+}, x_{-}) = (\alpha^{-1}, N_{c})$ as a strongly motivated
conjecture.

This note is the math--physics bridge: it draws structural consequences
of the $\chi_{-4}$-unification of \cite{paper-gstar-A} for the FTD
identification, without proving the identification itself. All
load-bearing claims are tagged \texttt{[SYNTHESIS]} or \texttt{[STRONGLY
MOTIVATED CONJECTURE]} as appropriate; no new mathematical theorem is
claimed in this note.

\section{The value-set coincidence and the Gaussian-integer-module ontology}\label{sec:ontology}

\begin{observation}[Value-set coincidence]\label{obs:value-set-coincidence}
\textbf{[Synthesis.]} Three independent starting points terminate at the
same three-element set:
\begin{enumerate}[label=\textup{(\arabic*)}, leftmargin=*, itemsep=2pt]
\item the FTD physical primitive --- each lattice site $v$ takes a state
$s_v \in \{-1, 0, +1\}$ encoding \emph{negative state}, \emph{void},
\emph{positive state} \cite{ftd-spec};
\item the value set of the unique non-trivial real Dirichlet character of
conductor 4 \cite[\S 16]{paper-gstar-A};
\item the closure of the imaginary-unit orbit $\{i^{n} : n \in \Z\}$
under the imaginary-part map: $\{\mathrm{Im}(i^{0}), \mathrm{Im}(i^{1}),
\mathrm{Im}(i^{2}), \mathrm{Im}(i^{3})\} = \{0, 1, 0, -1\}$.
\end{enumerate}
That these three sets coincide is a set-level identity, not yet a
\emph{functorial} identification. Whether the FTD voxel alphabet
\emph{arises from} the character $\chi_{-4}$, or only happens to share
its value set, is a physical-modelling question that lies outside the
scope of the mathematical paper.
\end{observation}

\begin{remark}[Candidate ontology: Z[i]-substrate realism]\label{rem:zi-substrate-realism}
\textbf{[Synthesis.]} The value-set coincidence of
Observation~\ref{obs:value-set-coincidence} suggests a working ontology
in which the empirical universe is the real-coordinate readout of a
$\Z[i]$-module structure. Layering the established FTD framework against
this ontology:
\begin{itemize}[leftmargin=*, itemsep=2pt]
\item voxels are positions in a $\Z[i]^d$-style lattice ($d = 3$ in current FTD);
\item voxel states $s_v \in \{-1, 0, +1\}$ are the values
$\mathrm{Im}(i^{\bullet})$ readouts of the unit group $\Z[i]^\times$;
\item the flux field $J \in \R^3$ is the continuous real part of an
underlying complex object;
\item time is the order parameter of the cyclic action $\langle i \rangle$
at long wavelength;
\item the mass--length calibration $a_{\mathrm{phys}} \equiv \ell_P$ is
the choice of unit on the $\Z[i]$-module;
\item dimensionless physical constants ($\alpha$, $N_c$, mass ratios) are
intrinsic invariants of the $\Z[i]$-module structure.
\end{itemize}
The ontological position is most accurately described as
\emph{Eleatic structural realism, specialized to the Gaussian-integer
module}: Eleatic because the substrate is a single algebraic object;
structural realist in the sense of Ladyman--Ross \cite{ladyman-ross},
where the entities are nodes in the structural network rather than
independent objects; specialized because the particular structure is the
maximal order of $\Q(i)$ rather than an arbitrary mathematical structure.
This is a sharper position than Tegmark's broad mathematical-universe
hypothesis \cite{tegmark2008}: not "all mathematical structures are
physical" but "the physical is the $\Z[i]$-module slice of the
mathematical." Wigner's "unreasonable effectiveness" \cite{wigner1960}
becomes, in this picture, eminently reasonable.
\end{remark}

\begin{remark}[Falsifiability]\label{rem:falsifiability}
\textbf{[Synthesis.]} The $\Z[i]$-substrate ontology is empirically
refutable: if experiment ever produced an observable that requires a
$\Z[\rho]$-module structure (Eisenstein, sixfold) that cannot be
embedded in a $\Z[i]$ ambient, or that requires
$\Z[\sqrt{-2}]$ or $\Z[\sqrt{-7}]$ (other class-number-1 quadratic
orders), the framework breaks. The position has therefore the
falsifiability profile of a substantive empirical hypothesis, not a
metaphysical claim.
\end{remark}

\section{The Born rule and $\chi_{-4}$ as the same arrow}\label{sec:born-rule-arrow}

The deepest single structural connection that the $\chi_{-4}$-unification
of \cite{paper-gstar-A} makes available is the identification of two
distinct operations as instances of the same arrow.

\begin{observation}[Born rule and $\chi_{-4}$ as the same arrow]\label{obs:born-arrow}
\textbf{[Synthesis.]} The Born rule of quantum mechanics and the
character $\chi_{-4}$ are formally the same projection applied to
different complex objects:
\begin{itemize}[leftmargin=*, itemsep=4pt]
\item \textbf{Born rule}: $\psi \in \C \;\mapsto\; |\psi|^{2} \in \R_{\geq 0}$,
projecting a complex amplitude onto its real modulus.
\item \textbf{Character $\chi_{-4}$}: $i^{n} \in \langle i\rangle \subset \C \;\mapsto\; \mathrm{Im}(i^{n}) \in \{-1, 0, +1\}$,
projecting a complex unit onto its signed imaginary coordinate.
\end{itemize}
Both are functions of the form
$(\text{complex structure}) \xrightarrow{(\text{single real coordinate})} \R$.
The Born rule reads the modulus; $\chi_{-4}$ reads the imaginary part.
They are dual projections from the same complex structure, both producing
ternary or near-ternary outcomes precisely because the imaginary-part
projection of the cyclic group $\langle i \rangle$ takes exactly three
distinct values.
\end{observation}

\begin{remark}[Reflexive projection at three scopes]\label{rem:three-scopes}
\textbf{[Synthesis.]} Under the FTD framework's identification of the
voxel state field with the substrate's reflexive projection
\cite{ftd-spec}, the same arrow applies at three
nested scopes:
\[
\underbrace{\text{reflexive projection of substrate}}_{\text{FTD voxel state} \;=\; \mathrm{Im}(i^{\bullet})}
\;\;\sim\;\;
\underbrace{\text{measurement event of amplitude}}_{\text{Born rule outcome} \;=\; |\psi|^{2}}
\;\;\sim\;\;
\underbrace{\text{imaginary-part trace of unit orbit}}_{\chi_{-4}(n) = \mathrm{Im}(i^{n})}.
\]
The first and third are mathematically equivalent (the FTD voxel state
field, viewed as a function on lattice sites with values in $\Z[i]^\times \cup \{0\}$
read off through the imaginary-part map, is literally the character
$\chi_{-4}$ on the angular index). The second is conjectured-structurally
equivalent to the first via the FTD reflexivity ontology; this is the
load-bearing physical conjecture and is tagged
\texttt{[STRONGLY MOTIVATED CONJECTURE]} accordingly.
\end{remark}

\section{The master quadratic as the bridge}\label{sec:master-quadratic-bridge}

The FTD bridge from $\chi_{-4}$-derived mathematics to physical observables
runs through the master quadratic $P_{\Gs}$. We recall (from
\cite[Theorem 16.2]{paper-gstar-A}) that $\chi_{-4}$ generates the
integer coefficient $16 = |\Z[i]^\times|^{2}$ (lattice projection) and
the constant $\Gs$ (Chowla--Selberg projection), but does \emph{not}
generate the specific polynomial form $x^{2} - 16\Gss x + 16\Gsss$ from
classical CM/Weber/Hecke machinery (the three negative tests of
\cite[\S 16.5]{paper-gstar-A}).

The exponent pair $(2, 3)$ on $\Gs$ in the polynomial coefficients is
nevertheless \emph{uniquely picked out} by joint root-matching:

\begin{observation}[Joint-matching uniqueness, restated from {\cite[\S 16.6]{paper-gstar-A}}]\label{obs:joint-matching-bridge}
\textbf{[Synthesis.]} Among polynomial forms $y^{2} - 16 R^{p} y + 16 R^{q}$
at $R = \Gs$ with $(p, q) \in \Z^{2}$, $|p|, |q| \leq 5$, only the pair
$(p, q) = (2, 3)$ produces an ordered root pair $(y_{+}, y_{-})$ that
simultaneously matches the FTD pair $(\alpha^{-1}, N_{c})$ to sub-percent
precision. The match precision at $(2, 3)$ is $1.26$~ppm for $y_{+} =
137.036$ and $0.8\%$ for $y_{-} = 3.024$. Other exponent pairs match
one root but not both, or neither.
\end{observation}

\begin{remark}[Information-theoretic minimality of the FTD bridge]\label{rem:bridge-minimality}
\textbf{[Synthesis.]} The FTD math--physics bridge requires exactly
\emph{one bit} of input from the physical side --- the exponent-pair
selector $(p, q) \in \{(2, 3), (3, 4), (4, 5), \ldots, (1, 2), (2, 4), \ldots\}$ ---
and that bit is over-determined by joint matching to two distinct
physical constants $(\alpha^{-1}, N_{c})$. Equivalently: the FTD master
quadratic is the unique polynomial of the form $y^{2} - 16 R^{p} y + 16 R^{q}$
at $R = \Gs$ whose ordered root pair $(y_{+}, y_{-})$ jointly matches
$(\alpha^{-1}, N_{c})$ to sub-percent precision.

This is not a derivation of the bridge --- the joint match itself is the
input that selects $(p, q)$ --- but it is the cleanest available
\emph{minimality statement}: the math--physics bridge in the FTD
framework requires no parameters beyond the choice of one exponent
pair, and the choice is over-determined (hence robust) against the joint
match to two distinct physical constants.
\end{remark}

\begin{observation}[The locus of the math--physics gap]\label{obs:gap-locus}
Combining \cite[\S 16.5]{paper-gstar-A} (no CM-internal source for $(2, 3)$)
with Observation~\ref{obs:joint-matching-bridge} (joint match uniquely
fixes $(2, 3)$): the gap between $\chi_{-4}$-derived mathematical
structure and the FTD physical identification is localised precisely at
the polynomial \emph{assembly} step. The integer coefficient $16$ and
the constant $\Gs$ are both forced by $\chi_{-4}$. The exponent pair
$(2, 3)$ is forced only by the joint physical-matching requirement.
Closing this last gap --- if a clean derivation of $(2, 3)$ from
$\chi_{-4}$-internal data exists --- would convert the FTD bridge from
a two-arrow conjecture into a single-arrow theorem.
\end{observation}

\section{The full descent chain}\label{sec:descent}

The cumulative content of \cite{paper-gstar-A} and this companion can be
displayed as a single descent chain from physics to ontological bedrock:
\[
\begin{array}{c}
\underbrace{(\alpha^{-1},\, N_{c})}_{\text{measured physical constants}} \\
\Big\uparrow \text{\quad\texttt{[STRONGLY MOTIVATED CONJECTURE]}\quad (FTD bridge)} \\
\underbrace{P_{\Gs}(x) = x^{2} - 16\Gss x + 16\Gsss}_{\text{master quadratic}} \\
\Big\uparrow \text{\quad(2, 3) exponents: uniquely forced by joint root match \cite[\S 16.6]{paper-gstar-A}} \\
\underbrace{\{16,\; \Gs\}}_{\text{building blocks}} \\
\Big\uparrow \text{\quad four functorial projections \cite[Thm 16.2]{paper-gstar-A}} \\
\underbrace{\chi_{-4}}_{\text{Kronecker character of } \Q(i)} \\
\Big\uparrow \text{\quad imaginary-part trace}\;\; \chi_{-4}(n) = \mathrm{Im}(i^{n}) = \sin(\pi n/2) \\
\underbrace{\{-1,\, 0,\, +1\}}_{\text{value set} \;=\; \text{FTD voxel alphabet (set-level)}} \\
\Big\uparrow \text{\quad closure of }\langle i\rangle\text{-orbit under squaring \& modulus} \\
\underbrace{i^{2} = -1}_{\text{ontological zero-point}}
\end{array}
\]
Every arrow above the FTD bridge is established (theorem, observation,
or honest conjecture) in \cite{paper-gstar-A}. The FTD bridge itself is
the physical conjecture that this companion records under
\texttt{[STRONGLY MOTIVATED CONJECTURE]}. Everything below the bridge is
pure mathematics; the only place that requires physical input is the
single bit of the exponent-pair selector.

\section{Outlook}\label{sec:outlook}

Two productive lines of work are identified by the descent chain.

\textbf{Closing the bridge.} If a derivation of the exponent pair
$(2, 3)$ from $\chi_{-4}$-internal data can be made rigorous --- the
candidate framing in \cite[\S 16.5, Conjecture (Sym\textsuperscript{2}--Sym\textsuperscript{3})]{paper-gstar-A}
is the most natural starting point --- the FTD bridge would become a
single-arrow theorem rather than a two-arrow conjecture. This would
not prove the empirical match to $(\alpha^{-1}, N_{c})$ (that remains an
experimental fact about the universe), but it would localise the
conjectural content to a single physical assertion: that the universe's
electroweak and strong sectors are built from a $\Z[i]$-substrate.

\textbf{Extending the bridge.} The current FTD identifications cover
$\alpha$ and $N_{c}$. The Standard Model has additional dimensionless
constants ($\sin^{2}\theta_{W}$, $m_{\mu}/m_{e}$, $m_{\tau}/m_{e}$,
PMNS angles, etc.) which the FTD framework currently parameterises
without first-principles derivation. The $\chi_{-4}$-unification suggests
that further dimensionless invariants may be discoverable as roots of
higher-rank master polynomials living in $\mathrm{Sym}^{\bullet}(H^{1}(E_{\mathrm{lemn}}))$,
or as Hecke L-values of related CM elliptic curves. This is the natural
program for extending the bridge from $(\alpha^{-1}, N_{c})$ to the full SM.

A second-tier program suggested by the polymath synthesis of
\cite[appendix]{paper-gstar-A} is the identification of the
$\chi_{-3}$-side (equianharmonic, $|\Aut| = 6$) as the structural home
of the QCD gauge structure, with $\chi_{-4}$ (lemniscatic) as the home
of the electroweak structure. The Moore Layer Theorem of FTD
\cite{ftd-spec} produces $U(1) \times SU(2) \times SU(3)$ from a
27-block polyhedral decomposition; whether this can be reformulated as
the joint action of $\chi_{-4}$ (giving $U(1)$ and part of $SU(2)$) and
$\chi_{-3}$ (giving $SU(3)$ and the rest of $SU(2)$ via the
Hurwitz-quaternion intersection $\mathbb{H} \supset \Z[i], \Z[\rho]$)
is an open speculative direction worth investigating but not pursued
here.

\begin{thebibliography}{9}

\bibitem{paper-gstar-A}
C. Pacitanimulli,
\textit{The Kronecker character $\chi_{-4}$ as the joint source of the
lemniscatic identity algebra},
manuscript, 2026.

\bibitem{ftd-spec}
C. Pacitanimulli,
\textit{Foundational Ternary Dynamics: specification and algebraic spine},
project documentation, 2026.
Available at the project repository; companion manuscript in preparation.

\bibitem{tegmark2008}
M.\ Tegmark,
\textit{The Mathematical Universe},
Foundations of Physics \textbf{38} (2008), 101--150.

\bibitem{ladyman-ross}
J.\ Ladyman, D.\ Ross, D.\ Spurrett, and J.\ Collier,
\emph{Every Thing Must Go: Metaphysics Naturalized},
Oxford University Press, 2007.

\bibitem{wigner1960}
E.\ P.\ Wigner,
\textit{The Unreasonable Effectiveness of Mathematics in the Natural
Sciences},
Communications on Pure and Applied Mathematics \textbf{13} (1960), 1--14.

\bibitem{chowlaselberg}
S.\ Chowla and A.\ Selberg,
\textit{On Epstein's zeta-function},
Journal f\"ur die reine und angewandte Mathematik \textbf{227} (1967),
86--110.

\bibitem{chudnovsky1976}
G.~V.\ Chudnovsky,
\textit{Algebraic independence of constants connected with the
exponential and elliptic functions},
Doklady Akademii Nauk SSSR \textbf{4} (1976), 698--701.

\bibitem{coates-wiles}
J.\ Coates and A.\ Wiles,
\textit{On the conjecture of Birch and Swinnerton-Dyer},
Inventiones Mathematicae \textbf{39} (1977), 223--251.

\end{thebibliography}

\end{document}
